\documentclass[11pt]{article}

\usepackage{fullpage}
\usepackage{amsmath}
\usepackage{epsfig}
\usepackage{graphics}
\usepackage{caption}
\usepackage{circuitikz}

\oddsidemargin -0.15in
\evensidemargin -0.15in
\headheight 0in
\headsep 0in
%\topmargin -0.25in
\textheight 9.3in
\textwidth 6.8in

\usepackage{enumitem,amssymb}
\newlist{todolist}{itemize}{2}
\setlist[todolist]{label=$\square$}

\begin{document}


\begin{center}
\section*{
Franklin \& Marshall College \hfill PHY 223\\
\bigskip
{\Large \it BB Scattering: Writing  (Part 2)}}
\end{center}

% \vspace{6mm}

\noindent {\Large {\bf Feedback on figures}}

\vspace{3mm}

You will be assigned 3 figures from another team. Follow the instructions on Canvas to submit comments that will help the other team improve their figures. \textbf{Make sure to do this before continuing!}

% For each of the figures, you need to comment on 3 features (shape, color, data displayed, concept represented, text, axes, etc.). Your comments can highlight things that you think are successful and/or a way to improve them. Keep the following questions in mind:\vspace{3mm}


% \begin{itemize}
%  \item Is the figure free of errors?
%  \item Is the figure visually appealing?
%  \item Are there elements of the figure that could be misleading or lead to misunderstandings?
%  \item Are there colors or font sizes that are hard to read or interpret?
%  \item Are the axes, ticks, grids, legends, etc. useful for communicating the results, or are they distracting?
%  \item Is there redundant information that could be simplified without loss of clarity?
%  \item Does the caption help connect the figure to the goals of the report?
%  \item Are there other elements from the style guide and other visualization tips that would be useful for the creators to consider?
% \end{itemize}

% \noindent Make sure to leave comments on \textbf{all three figures}. You don't have to address every question above, but consider each of them and make comments if appropriate. You do not need to include your names on your submission, but make sure that your comments are constructive and will be helpful to the readers---your instructor will read them as well. Then submit your comments as a PDF using the corresponding assigment on Canvas.


\vspace{5mm}

\noindent {\Large {\bf Mapping out your report}}

\vspace{3mm}

Before you start to write in earnest, you should plan the document with your teammates. This will make writing easier and help the text be more coherent.
\vspace{3mm}

\begin{enumerate}
 \item Start a document for your report. It doesn't matter what you use at this point as long as all the team members have concurrent access to it (Google Docs, Overleaf, etc.).

 \item Create section titles for your report:

  \begin{itemize}
   \item Title page including an abstract
   \item Introduction
   \item Theory
   \item Methods
   \item Data and analysis
   \item Conclusion
   \item Team Roles and Acknowledgements
   \item References
  \end{itemize}

 \item Insert your figures to the appropriate sections. Note that you may wish to replace your current versions of the figures once you receive feedback on them, but you can still put them in place to plan around them.

%  \item In each section, write a bullet point list of the main points you want to make. These ideas can definitely change as you write your report and your thoughts become clearer! Make sure that the narrative of your results across the whole report is represented in your outline. \textit{Note: your final report will not include these bulleted lists---as you write the full version of your report, you will replace the bullets with paragraphs.}

\item Review the instructions for each of the sections above in the ``Guide to writing a lab report'' linked to Canvas. Then go through each section of your lab report and create some bullet points identifying the main things you want to include and the main points you want to make. These ideas can definitely change as you write your report and your thoughts become clearer! Make sure that the narrative of your results across the whole report is represented in your outline. \textit{Note: your final report will not include these bulleted lists---as you write the full version of your report, you will replace the bullets with paragraphs.}

\end{enumerate}

\vspace{1mm}

\noindent Once you have some notes in each section, have each member of your team choose a section to start filling in more detail. Some ideas for getting started are given below:

\vspace{5mm}

% \newpage

\noindent {\Large {\bf Starting your Introduction}}

\vspace{3mm}

Review the section in the ``Guide to write a lab report'' about introductions. If you have not already done so, do some research on about scattering experiments in physics (see instructions in the previous writing assignment). Add bullets or sentences summarizing some of these experiments and how our BB scattering experiment relates to the physics of those other experiments. At the same time, add references to the external sources where you found this information in your References section.

% With your teammates, write a draft introduction in your document. Make sure to use information from your research about different types of scattering experiments, and include references to existing work.

\noindent {\Large {\bf Starting your Methods}}

\vspace{3mm}

Review the section in the ``Guide to write a lab report'' about Methods. Write down the steps by which you collected the data, including references to your experimental setup figure. Write as if the reader had not already done the experiment themselves and did not have access to the instructions that you followed. Also make sure to include any specific choices you made that might the precision and accuracy of your result---you may wish to review the Data Acquisition assignment you completed where you described any such choices. You may also reference the instructions provided for the data acquisition, but make sure to describe your methods \textit{in your own words!}

% With your teammates, write a draft introduction in your document. Make sure to use information from your research about different types of scattering experiments, and include references to existing work.

\vspace{5mm}

\noindent {\Large {\bf Starting your Theory}}

\vspace{3mm}

Review the section in the ``Guide to write a lab report'' about Theory. Motivate the mathematical expressions you will use for your analysis and define the variables you will be using. If your have a figure that relates to this section, explicitly describe how the equations are related to the figure and the experiment. You should have math here!

\vspace{5mm}

\noindent {\Large {\bf Starting your Data and Analysis}}

\vspace{3mm}

Review the section in the ``Guide to write a lab report'' about Data and Analysis. Describe your data and analysis figures, including the manner in which you processed the data raw data into your final figure(s), their uncertainties, and how you performed any fits.  Describe your final radius measurement (including uncertainty) and compare it to a physical measurement of the radius, and interpret the agreement or disagreement between those values. 

% With your teammates, write a draft introduction in your document. Make sure to use information from your research about different types of scattering experiments, and include references to existing work.

\vspace{5mm}

\noindent {\Large {\bf Planning the rest of the work}}

\vspace{3mm}

Before leaving for the day, make a plan with your colleagues about who will do what with the rest of the report and when that work is due. You should leave time for edits when the document is fully written. Many lab reports read as a collection of text written by different people, so try to avoid this pitfall.
\vspace{6mm}

\noindent {\bf When you are done for today, upload your outline (and any additional text you have filled in) to Canvas as a PDF.}

\end{document}
